\scanpage{167}
\begin{center}- 45 -\end{center}\hfill IX

\uline{correspondantes sont respectivement} $(q^n\epsilon_i)_{1\leq i\leq\mu}$ \uline{et} $(\epsilon_i)_{1\leq i\leq\mu}$, \uline{où les} $\epsilon_i$ $(1\leq i\leq\mu)$ \uline{sont des} racines de l'unité.

En fait, la formule (4.1.2) donne pour le dernier morceau les valeurs propres $q^n/(q^n\epsilon_i)=\epsilon_i^{-1}$, mais comme $\epsilon_i$ est une racine de l'unité, on a $\epsilon_i^{-1}=\overline{\epsilon_i}$ (conjugué complexe), et comme la famille des $\epsilon_i$ est stable par conjugaison complexe (comme famille des solutions d'une équation à coefficients entiers), la famille des $\epsilon_i^{-1}$ est aussi celle des $\epsilon_i$, ce qui prouve la dernière assertion de 4.4.

4.5. Nous allons donner une application de 4.3 a) à la théorie du conducteur local des schémas abéliens. Supposons $k$ algébriquement clos (*). Rappelons (SGA 5 X 6.2 ou [20])que pour tout schéma en $\ell$-groupes commutatifs étale $F_K$ sur $K$, $(\ell\ne p)$, on définit un entier
\[
 (4.5.1)\qquad\alpha_s(F_K)=\alpha_s(i_*(F_K))
\]
($i:\eta\longrightarrow S$ désignant l'inclusion), de la façon suivante : on choisit une extension galoisienne finie $K'$ de $K$, de groupe $\Gamma$, qui déploie $F_K$, de sorte que $F_K$ soit défini par un $\ell$-groupe fini $M$ sur lequel $\Gamma$ opère.
D'autre part, on associe à l'extension $K'$ de $K$ un certain $\mathbb Z_\ell[\Gamma]$-\uline{module projectif} (défini à isomorphisme près)
\[
 (4.5.2)\qquad\mathrm{Sw}=\mathrm{Sw}(K'/K,\mathbb Z_\ell)\ ,
\]
et on pose
\[
 (4.5.3)\qquad\alpha_s(F_K)=\operatorname{long}\operatorname{Hom}_{\mathbb Z_\ell(\Gamma)}(\mathrm{Sw},M)\ .
\]
On vérifie (loc. cit) que cette expression ne dépend pas du choix de l'extension $K'/K$. D'ailleurs, le fait que $\mathrm{Sw}$ soit un $\mathbb Z_\ell[\Gamma]$-module
\figure{\begin{tikzpicture}\draw (0,0)--(3,0);\end{tikzpicture}}
(*) Il suffirait en fait que $k$ soit \uline{parfait}, mais par passage au hensélisé strict de $S$, on se ramène aussitôt au cas envisagé.
\margin{\unsure{263|363}}
\margin{363}
\margin{166}
